% Add \usepackage{booktabs} to your document preamble.
% Use table* for a two-column paper; use table for a single-column document.
\begin{table*}[t]
  \centering
  \caption{Batch query latency and byte cache hit ratio. CELEBI combines page
    caching, query reordering, and future-aware eviction. Latency reduction is
    relative to the baseline.}
  \label{tab:batch-latency-cache-hits}
  \small
  \setlength{\tabcolsep}{6pt}
  \begin{tabular}{@{}lrrrrrr@{}}
    \toprule
    Workload & \multicolumn{3}{c}{Batch query latency (s)}
      & \multicolumn{2}{c}{Byte cache hit ratio (\%)}
      & \shortstack{CELEBI latency\\reduction (\%)} \\
    \cmidrule(lr){2-4}\cmidrule(lr){5-6}
      & Baseline & Paging & CELEBI & Paging & CELEBI & \\
    \midrule
    ClickBench (43 queries)
      & 43.17 & 39.34 & \textbf{37.84}
      & 22.53 & \textbf{37.01} & \textbf{12.35} \\
    TPC-H SF50 (22 queries)
      & 32.56 & 30.75 & \textbf{23.82}
      & 9.01 & \textbf{36.79} & \textbf{26.86} \\
    \bottomrule
  \end{tabular}
  \par\vspace{4pt}
  \begin{minipage}{\textwidth}
    \footnotesize
    Batch latency is the sum of per-query \texttt{execute().fetchall()} elapsed
    times, averaged over executions 2 and 3; execution 1 is excluded as warmup.
    It excludes query-label and cache-hint calls and reorder preparation.
    Byte cache hit ratio is the sum of cache-served decoded value bytes divided
    by the sum of requested decoded value bytes, pooled over executions 2 and 3.
    Repeated requests are included; string offsets, null masks, and allocation
    padding are excluded. Latency and cache hits are collected in separate
    trace-off and trace-on runs with matching settings.
    ClickBench includes q28 and q29 and was remeasured after the AVG optimization;
    TPC-H uses the previous measurements. Paging uses LRU without reordering.
  \end{minipage}
\end{table*}
